% Part: first-order-logic
% Chapter: syntax-and-semantics
% Section: unique-readability

\documentclass[../../../include/open-logic-section]{subfiles}

\begin{document}

\olfileid{fol}{syn}{unq}

\olsection{اکّو طرح نال پڑھت}

\begin{explain}
جس طرح اسیں !!{formula}s دی تعریف کیتی اے، اوہ ضمانت
دیندا اے کہ ہر !!{formula} دی \emph{اکّو پڑھت} ہووے؛
یعنی !!{formula}s دے تشکیل قاعدیاں مطابق اوہنوں بناؤن
دا بنیادی طور تے صرف اک طریقہ اے، تے اوہدی نحوی ”تعبیر“
دا وی صرف اک طریقہ اے۔ جے ایہ نہ ہوندا، تاں مبہم
!!{formula}s ہوندے، یعنی اوہ !!{formula}s جیہناں
دیاں اک توں ودھ پڑھتاں یا نحوی تعبیراں ہوندیاں—ایہ صاف
اے کہ اسیں اس توں بچنا چاہندے آں۔ پر اس توں وی
اہم گل ایہ اے کہ ایس خاصیت دے بغیر اگلیاں
زیادہ تر تعریفاں تے ثبوت کم نہیں کرن گے۔

ایہ گل واضح کرن دا شاید سب توں چنگا طریقہ ایہ ویکھنا
اے کہ جے اسیں !!{formula}s دے تشکیل لئی خراب قاعدے
دتے ہوندے، جیہڑے اکّو طرح نال پڑھت دی ضمانت نہ
دیندے، تاں کی ہوندا۔ مثال لئی، اسیں رابطیاں دے
تشکیل قاعدیاں وچ قوساں بھل سکدے ساں؛ مثلاً اسیں
ایہ اجازت دے سکدے ساں:
\begin{quote}
جے $!A$ تے $!B$، !!{formula}s نیں، تاں $!A \lif !B$ وی اے۔
\end{quote}
ایٹمی فارمولا $!D$ توں شروع کر کے ایہ سانوں $!D
\lif !D$ بناؤن دی اجازت دیندا۔ ایس توں، $!D$ دے نال،
سانوں $!D \lif !D
\lif !D$ مل جاندا۔ پر ایہ کرن دے دو طریقے نیں:
\begin{enumerate}
\item اسیں $!D$ نوں $!A$ منّیئے تے $!D \lif !D$ نوں $!B$۔
\item اسیں $!A$ نوں $!D \lif !D$ منّیئے تے $!B$،~$!D$ اے۔
\end{enumerate}
ایس دے مطابق !!{formula}~$!D \lif !D \lif !D$
نوں ”پڑھن“ دے دو طریقے نیں۔ ایہ $!B
\lif !C$ دے روپ وچ اے، جدوں $!B$، $!D$ اے تے $!C$،
$!D \lif !D$ اے؛ پر \emph{ایہ ایس روپ وچ وی اے}:
$!B \lif !C$، جتھے $!B$، $!D \lif !D$ اے تے $!C$
اوہو~$!D$ اے۔

جے ایہ ہووے، تاں ساڈیاں تعریفاں ہر وار کم نہیں کرن
گیاں۔ مثال لئی، جدوں اسیں کسے فارمولے دے
!!{main operator} دی تعریف کردے آں، تاں آکھدے آں:
$!B \lif !C$ روپ دے فارمولے وچ !!{main operator}،
~$\lif$ دی دکھائی گئی آمد اے۔ پر جے اسیں فارمولا $!D
\lif !D \lif !D$ نوں اُتے دسے دو وکھرے طریقیاں نال
$!B \lif !C$ دے روپ وچ رکھ سکدے آں، تاں پہلی
صورت وچ $\lif$ دی پہلی آمد !!{main operator} بنے گی
تے دوجی صورت وچ دوجی آمد۔ پر ساڈا مقصد
!!{main operator} نوں !!{formula} دا اک \emph{فنکشن}
% PNB-SYN-006: اگلی main-operator تعریف مطابق ایٹمی فارمولیاں دا مرکزی عامل نہیں؛ضمنی غیرایٹمی دائرہ صاف۔
منّنا اے؛ یعنی ہر غیر ایٹمی !!{formula} وچ !!{main
  operator} دی ٹھیک اک آمد ہونی چاہیدی اے۔
\end{explain}

\begin{lem}
!!a{formula}~$!A$ وچ کھبیاں تے سجیاں قوساں دی
گنتی برابر اے۔
\end{lem}

\begin{proof}
اسیں ایہ $!A$ دی تشکیل اُتے استقرا نال ثابت کردے آں۔
ایس لئی دو گلاں چاہیدیاں نیں: (الف) پہلاں ثابت کریے
کہ سارے ایٹمی فارمولیاں وچ متعلقہ خاصیت اے
(استقرا دی بنیاد)۔ (ب) فیر ثابت کریے کہ جدوں دتے
فارمولیاں توں نویں فارمولے بنائے جان، تاں جے پرانیاں
وچ خاصیت اے، نویاں وچ وی ہووے۔

$l(!A)$ نوں کھبیاں قوساں دی گنتی تے $r(!A)$ نوں
~$!A$ وچ سجیاں قوساں دی گنتی منّو؛ ایس طرح
$l(t)$ تے $r(t)$، ٹرم~$t$ وچ کھبیاں تے سجیاں
قوساں دی گنتی نیں۔

\begin{prob}
    ہر ٹرم~$t$ لئی $l(t) = r(t)$ ثابت کرو۔
\end{prob}

\begin{enumerate}
\tagitem{prvFalse}{\indcase{!A}{\lfalse}{$\indfrm$ وچ $0$ کھبیاں تے $0$
    سجیاں قوساں نیں۔}}{}

\tagitem{prvTrue}{\indcase{!A}{\ltrue}{$\indfrm$ وچ $0$ کھبیاں تے $0$
    سجیاں قوساں نیں۔}}{}

\item \indcase{!A}{\Atom{R}{t_1,\dots,t_n}}{$l(\indfrm) = 1 + l(t_1) +
  \dots + l(t_n) = 1 + r(t_1) + \dots + r(t_n) = r(\indfrm)$۔ ایتھے
  اسیں اوہ گل ورتدے آں جیہڑی مشق لئی چھڈی گئی اے:
  $l(t) = r(t)$، ہر ٹرم~$t$ لئی۔}

\item \indcase{!A}{\eq[t_1][t_2]}{$l(\indfrm) = l(t_1) + l(t_2) =
  r(t_1) + r(t_2) = r(\indfrm)$۔}

\tagitem{prvNot}{\indcase{!A}{\lnot !B}{استقرائی مفروضے نال
    $l(!B) = r(!B)$۔ اس لئی $l(\indfrm) = l(!B) = r(!B) =
    r(\indfrm)$۔}}{}

\item \indcase{!A}{(!B \ast !C)}{استقرائی مفروضے نال $l(!B) =
  r(!B)$ تے $l(!C) = r(!C)$۔ اس لئی $l(\indfrm) = 1 + l(!B) + l(!C) =
  1 + r(!B) + r(!C) = r(\indfrm)$۔}

\tagitem{prvAll}{\indcase{!A}{\lforall[x][!B]}{استقرائی
    مفروضے نال $l(!B) = r(!B)$۔ اس لئی $l(\indfrm) = l(!B) = r(!B) =
    r(\indfrm)$۔}}{}

% Print case for \lexists if prvEx and not prvAll
\tagitem{prvAll,notprvEx}{}{\indcase{!A}{\lexists[x][!B]}{استقرائی
    مفروضے نال $l(!B) = r(!B)$۔ اس لئی $l(\indfrm) = l(!B) = r(!B) =
    r(\indfrm)$۔}}

% Just say similarly if prvEx and prvAll
\tagitem{notprvAll,notprvEx}{}{\indcase{!A}{\lexists[x][!B]}{ایویں ای۔}}
\end{enumerate}
\end{proof}

\begin{defn}[پورے توں چھوٹا مُڈھلا ٹکڑا]
علامتاں دی لڑی $!B$، علامتاں دی لڑی~$!A$ دا
\emph{پورے توں چھوٹا مُڈھلا ٹکڑا} اے، جے $!B$
نال علامتاں دی اک غیر خالی لڑی جوڑ کے~$!A$ ملے۔
\end{defn}

\begin{lem}\ollabel{lem:no-prefix}
جے $!A$، !!a{formula} اے تے $!B$، $!A$ دا
پورے توں چھوٹا مُڈھلا ٹکڑا اے، تاں $!B$،
!!a{formula} نہیں اے۔
\end{lem}

\begin{proof}
مشق۔
\end{proof}

\begin{prob}
\olref[fol][syn][unq]{lem:no-prefix} ثابت کرو۔
\end{prob}

\begin{prop}
\ollabel{prop:unique-atomic}
جے $!A$ ایٹمی !!{formula} اے، تاں اوہ ہیٹھ
دیتیاں شرطاں وچوں ٹھیک اک شرط پوری کردا اے۔
\begin{enumerate}
\tagitem{prvFalse}{$!A \ident \lfalse$۔}{}
\tagitem{prvTrue}{$!A \ident \ltrue$۔}{}
\item $!A \ident \Atom{R}{t_1,\dots,t_n}$، جتھے $R$ اک $n$-مقامی
  !!{predicate} اے، $t_1$, \dots, $t_n$ ٹرم نیں، تے $R$،
  $t_1$, \dots, $t_n$ وچوں ہر اک اکّو طرح متعین اے۔
\item $!A \ident \eq[t_1][t_2]$، جتھے $t_1$ تے $t_2$
  اکّو طرح متعین ٹرم نیں۔
\end{enumerate}
\end{prop}

\begin{proof}
مشق۔
\end{proof}

\begin{prob}
\olref[fol][syn][unq]{prop:unique-atomic} ثابت کرو (اشارہ:
ٹرماں لئی \olref[fol][syn][unq]{lem:no-prefix} دا اک روپ
بیان کرو تے ثابت کرو۔)
\end{prob}

\begin{prop}[اکّو طرح نال پڑھت]
ہر !!{formula} ہیٹھ دیتیاں شرطاں وچوں ٹھیک اک
شرط پوری کردا اے۔
\begin{enumerate}
\item $!A$ ایٹمی اے۔

\tagitem{prvNot}{$!A$، $\lnot !B$ دے روپ وچ اے۔}{}

\tagitem{prvAnd}{$!A$، $(!B \land !C)$ دے روپ وچ اے۔}{}

\tagitem{prvOr}{$!A$، $(!B \lor !C)$ دے روپ وچ اے۔}{}

\tagitem{prvIf}{$!A$، $(!B \lif !C)$ دے روپ وچ اے۔}{}

\tagitem{prvIff}{$!A$، $(!B \liff !C)$ دے روپ وچ اے۔}{}

\tagitem{prvAll}{$!A$، $\lforall[x][!B]$ دے روپ وچ اے۔}{}

\tagitem{prvEx}{$!A$، $\lexists[x][!B]$ دے روپ وچ اے۔}{}
\end{enumerate}
ہور ایہ کہ ہر صورت وچ $!B$، یا $!B$ تے $!C$،
اکّو طرح متعین نیں۔ ایس دا مطلب اے کہ، مثال لئی،
وکھرے جوڑے $!B$, $!C$ تے $!B'$, $!C'$ نہیں ہو
سکدے جیہناں لئی $!A$ دونیں روپاں وچ ہووے:
\iftag{prvIf}{$(!B \lif !C)$ تے $(!B' \lif !C')$۔}{
    \iftag{prvOr}{$(!B \lor !C)$ تے $(!B' \lor !C')$۔}{
      \iftag{prvAnd}{$(!B \land !C)$ تے $(!B' \land !C')$۔}{}}}
\end{prop}

\begin{proof}
تشکیل دے قاعدے لازمی کردے نیں کہ جے !!a{formula}
ایٹمی نہیں، تاں اوہ کھلدی قوس~(،
\iftag{prvNot}{$\lnot$،}{} یا مسوّر نال شروع ہووے۔
دوجے پاسے، جیہڑا وی !!{formula} ہیٹھ دیاں
علامتاں وچوں کسے اک نال شروع ہووے، اوہ ایٹمی
ہووے گا: !!a{predicate}، !!a{function}، !!a{constant}\iftag{prvFalse}{، $\lfalse$}{}\iftag{prvTrue}{، $\ltrue$}{}۔

اس لئی سانوں اصل وچ صرف ایہ دکھانا اے کہ جے $!A$
اک پاسے $(!B \ast
!C)$ دے روپ وچ اے تے دوجے پاسے $(!B' \mathbin{\ast'} !C')$،
تاں $!B \ident
!B'$، $!C \ident !C'$، تے $\ast = {\ast'}$۔

فرض کرو $!A \ident (!B \ast !C)$ وی اے تے $!A \ident (!B'
\mathbin{\ast'} !C')$ وی۔ تاں یا $!B \ident !B'$، یا ایہ
نہیں۔ جے ایہ اے، تاں صاف اے کہ $\ast = {\ast'}$ تے
$!C \ident !C'$، کیوں جے ایہ فیر $!A$ دیاں
اوہ ذیلی لڑیاں نیں جیہڑیاں اکّو مقام توں شروع ہوندیاں
تے اکّو لمبائی رکھدیاں نیں۔ دوجی صورت $!B \not\ident !B'$
اے۔ کیوں جے $!B$ تے
$!B'$ دونیں $!A$ دیاں ذیلی لڑیاں نیں تے اکّو مقام
توں شروع ہوندیاں نیں، تاں اک دوجی دا پورے توں
چھوٹا مُڈھلا ٹکڑا لازماً ہووے گی۔ پر ایہ
\olref{lem:no-prefix} دے مطابق ناممکن اے۔
\end{proof}


\end{document}
